% ============================================================
% 1. INTRODUCTION -- humanized pass. Break symmetric signposting,
% vary constructions, keep the three contribution bullets (they work).
% ============================================================
\section{Introduction}

Anticipating what a person will do next from egocentric video is hard for a
structural reason: the most discriminative evidence, the action itself, has not
happened yet. The model has to work from preparatory cues instead. Human
behavioral signals are a natural candidate for those cues. Gaze leads
manipulation by something like half a second to a second, and the hands are the
effector that carries out intention, so conditioning a predictive model on gaze
and hand pose ought to help it anticipate.

V-JEPA~2 \citep{assran2025vjepa2} makes that idea testable. It learns a frozen
encoder and a latent predictor by predicting masked spatiotemporal features, and
its action-conditioned variant conditions the predictor on robot actions. We ask
whether the same predictor can be conditioned on \emph{human} signals instead,
gaze and hand pose from egocentric recordings, and whether that improves action
anticipation on EPIC-KITCHENS-100 (EK100) \citep{damen2022ek100}. Earlier phases
of the project had already returned a stubborn null at the two-second horizon,
and a control family had suggested that what looked like an effect of
conditioning was mostly an effect of feeding the model more tokens. That left one
question unanswered, with two answers that point in opposite directions. Perhaps
the injection pathway is too weak to expose whatever the signal carries; perhaps
the signal simply does not carry action-relevant information at this horizon. The
first answer argues for a stronger architecture. The second argues for giving up
on input-level conditioning altogether.

This paper sets out to decide between those two rather than assume either. The
contribution is diagnostic, and we state it plainly, null included:
\begin{itemize}
  \item We make the existing peer-token pathway a bit-exact control and guard
        against three silent-failure modes that can each fake a null, so that the
        null we report is a property of the signal and the pathway, not an
        artifact of the evaluation code (\S\ref{subsec:method-weak}).
  \item We run five diagnostics directly against the frozen predictor and the
        behavioral signal. They tell a consistent story: the signal reaches the output, the predictor
        ignores its content, and on its own the signal anticipates the next
        action barely above a class prior (\S\ref{subsec:diagnostics}). Varying
        the training rather than the input confirms it: a predictor whose
        training never included the signal matches one trained with it
        (\S\ref{subsec:training-ablation}). That evidence let us retire a stronger,
        GazeQwen-informed \citep{pham2026gazeqwen} pathway at the design stage
        instead of after building it (\S\ref{subsec:b4}).
  \item We isolate the one effect that does survive, a motion-specific advantage
        of behavioral signals on verbs over vision, and show that under a fair
        comparator it holds only in direction, never above the verb prior, and
        that the real constraint is single-participant sample size, not a floor
        in dynamic range (\S\ref{subsec:dissociation}--\ref{subsec:ddn}).
\end{itemize}
The result is a null we can explain rather than one we merely report, and the
explanation points past input-level conditioning toward representation-level
supervision of behavioral signals (\S\ref{sec:conclusion}). We keep every anticipation claim bounded to one participant at roughly one to two
seconds; the training-side ablations reach across participants, and we flag which
is which where it matters.